\section{Introducción}
\label{ch:introduccion}

El proyecto \textit{Semáforo de Insumos IMSS} se enfoca en el control, gestión y optimización del flujo de medicamentos e insumos dentro de unidades médicas del Instituto Mexicano del Seguro Social. El problema central no es la falta de información, sino su dispersión, el retraso con el que llega y la dificultad para convertirla en decisiones operativas claras. En la práctica se siguen procesos manuales, reportes tardíos y ajustes reactivos que terminan en desabasto, exceso de inventario o pérdida de trazabilidad.

Este documento corresponde al primer avance formal de la tesina. La investigación se encuentra en fase de campo y desarrollo, por lo que el contenido se concentra en el contexto, el problema, los objetivos, los antecedentes y el marco teórico tentativo. No se presentan resultados ni conclusiones porque todavía están en construcción.

\section{Objetivo general y específicos}
\label{sec:objetivos}

\subsection{Objetivo general}
\label{subsec:objetivo_general}

Diseñar e implementar un sistema de semaforización analítica para el control de insumos y medicamentos en unidades médicas del IMSS, capaz de capturar datos en tiempo real y transformar esa información en alertas visuales para identificar desabastos, excesos y cuellos de botella.

\subsection{Objetivos específicos}
\label{subsec:objetivos_especificos}

\begin{enumerate}
    \item Realizar el levantamiento de requerimientos en campo, considerando procesos reales de registro, surtido y consumo.
    \item Estructurar la base de datos que soporte el registro de entradas, salidas, niveles de stock y puntos críticos.
    \item Definir la lógica del algoritmo de semaforización, incluyendo reglas de negocio y umbrales por tipo de insumo.
    \item Diseñar una interfaz de usuario clara y operativa que permita visualizar el semáforo y facilitar la toma de decisiones.
\end{enumerate}

\section{Planteamiento del problema}
\label{sec:problema}

En las unidades médicas del IMSS se manejan grandes volúmenes de insumos con criticidad distinta. Sin embargo, el control del inventario sigue apoyándose en capturas manuales, reportes no estandarizados y validaciones tardías. Esta dinámica provoca que la información sobre el stock no sea confiable ni oportuna, lo que complica la planeación de compras, el surtido interno y la atención diaria a los pacientes.

El problema no es únicamente logístico; es analítico. Se requiere un mecanismo sencillo para convertir datos dispersos en señales claras: qué insumos están en riesgo, cuáles tienen exceso y cuáles están en un rango saludable. Sin esa lectura rápida y sistemática, las decisiones se toman tarde y con poca evidencia.

\section{Antecedentes}
\label{sec:antecedentes}

\subsection{Antecedentes de la empresa}
\label{subsec:antecedentes_empresa}

El Instituto Mexicano del Seguro Social es una de las instituciones de salud pública más grandes del país y opera una red amplia de hospitales y unidades médicas. Su operación exige disponibilidad continua de medicamentos, material de curación y consumibles clínicos. La gestión de inventarios en este contexto no solo impacta costos, también afecta de manera directa la calidad del servicio y la seguridad del paciente. Por ello, la eficiencia en el control de insumos es una prioridad operativa constante.

\subsection{Antecedentes teóricos}
\label{subsec:antecedentes_teoricos}

En la literatura y en la práctica hospitalaria se utilizan distintos enfoques para controlar inventarios: métodos de máximos y mínimos, puntos de reorden, clasificaciones ABC, reposición Just in Time y esquemas de inventario administrado por el proveedor (VMI). Estos modelos buscan equilibrio entre disponibilidad y costo, pero suelen depender de datos confiables y de un monitoreo continuo. En el sector salud, la variabilidad en la demanda y la criticidad de algunos insumos hacen que la simple regla de reorden no sea suficiente. En ese punto, un sistema de semaforización aporta una capa analítica que traduce datos en decisiones rápidas y visibles.

\subsubsection{Tabla comparativa}
\label{subsubsec:tabla_comparativa}

\begin{table}[!tbh]
    \centering
    \caption{Comparativa de enfoques de control de inventarios médicos.}
    \label{tab:comparativa_inventarios}
    \begin{tabular}{|L{3.0cm}|L{4.0cm}|L{4.0cm}|L{4.0cm}|}
        \hline
        \textbf{Sistema o metodología} & \textbf{Pros} & \textbf{Contras} & \textbf{Valor analítico de la propuesta} \\
        \hline
        Máximos--mínimos con punto de reorden & Fácil de implementar y auditar; estandariza el reabastecimiento. & No considera criticidad clínica ni variaciones diarias; depende de capturas puntuales. & El semáforo agrega prioridad por criticidad y señales tempranas ante tendencia de desabasto. \\
        \hline
        Just in Time / Kanban & Reduce sobreinventario y promueve flujo continuo. & Alta dependencia de proveedores; riesgo ante picos de demanda o contingencias. & La semaforización permite anticipar riesgos sin romper el flujo, con alertas operativas visibles. \\
        \hline
        VMI / consignación & Traslada parte de la gestión al proveedor; reduce carga administrativa interna. & Requiere datos confiables y acuerdos contractuales; menor control interno. & El sistema aporta trazabilidad interna y evidencia para negociar niveles y tiempos de reposición. \\
        \hline
    \end{tabular}
\end{table}

\section{Justificación}
\label{sec:justificacion}

Un sistema de semaforización analítica no pretende reemplazar los procesos existentes, sino hacerlos visibles y accionables. Al transformar datos de inventario en alertas claras, se reduce el tiempo de reacción ante desabastos, se evita el sobreabasto y se facilita la priorización de compras. Esto tiene impacto directo en la continuidad del servicio y en el uso responsable de recursos públicos.

Desde la perspectiva académica y profesional, el proyecto integra análisis de datos, diseño de bases de datos e interfaces centradas en la operación real. La propuesta es viable porque se apoya en datos que ya se generan en campo, pero que hoy no se explotan de manera sistemática.

\section{Alcances y limitaciones}
\label{sec:alcances_limitaciones}

\subsection{Alcances}
\label{subsec:alcances}

\begin{itemize}
    \item Definir el modelo de datos para registrar entradas, salidas, niveles de stock y umbrales operativos.
    \item Construir la lógica de semaforización para clasificar los insumos en verde, amarillo o rojo según su nivel y criticidad.
    \item Diseñar una interfaz de usuario funcional para la visualización del estado del inventario.
    \item Documentar el flujo de información desde la captura en campo hasta la visualización en el semáforo.
\end{itemize}

\subsection{Limitaciones}
\label{subsec:limitaciones}

\begin{itemize}
    \item El alcance inicial se enfoca en una unidad médica y en un subconjunto de insumos críticos, sujeto a disponibilidad de información.
    \item No se contempla, por ahora, la integración directa con sistemas institucionales externos ni con proveedores.
    \item La calidad de las alertas dependerá de la consistencia de los registros en campo.
    \item No se incluyen resultados ni mediciones de impacto, ya que el proyecto está en fase de desarrollo.
\end{itemize}

\section{Marco teórico tentativo}
\label{sec:marco_teorico_tentativo}

\begin{itemize}
    \item \textbf{Sistemas de gestión de inventarios en salud.} Conceptos de stock de seguridad, punto de reorden, clasificación ABC/XYZ, VMI y Just in Time aplicados a entornos clínicos.
    \item \textbf{Arquitectura de software para gestión de datos.} Modelado de bases de datos, integridad y normalización, flujo de datos transaccional y estructuras para trazabilidad.
    \item \textbf{Metodologías de desarrollo ágil.} Enfoques Scrum/Kanban para iteración rápida, validación con usuarios y priorización de requerimientos.
    \item \textbf{Herramientas de visualización de datos.} Tableros operativos, indicadores clave, semaforización y criterios de usabilidad para decisiones en tiempo real.
\end{itemize}
